

In this section, we prove Theorem~\ref{thm:smallsolutions}, which we
repeat here for convenience.

\medskip
\noindent
	\textbf{Theorem~\ref{thm:smallsolutions}} (\textbf{Bounds on solutions of linear
	systems}).
\hfill\\
	\emph{
Every satisfiable $d$-dim linear system $S$ 
 has a solution $x\in\setZ^d$ such that
	\[\nnorm{x}\leq
		\mlcm{S} \cdot \left(1+ (d+1)\,!\cdot (d\cdot
		\nnorm{S}+1)^{d}\right).\]
}
\medskip


We first prove the theorem for the case where the linear system $S$
	does not contain any divisibility constraints. We use the
	following result, which is a weaker form of the theorem from~\cite{Gathen1978}.

	\begin{lemC}[\cite{Gathen1978} (\textbf{Bounds on solutions to
		inequalities 
		via subdeterminants})]\label{lem:Gathen}
Let $C$ and $c$ be $n\times d$ and $n\times 1$ integer matrices, respectively,
and let $r$ be the rank of $C$. Let $m$ be an
	upper bound on the absolute values of the $(r-1)\times (r-1)$
	and $r\times r$ subdeterminants of the $n\times (d+1)$
	augmented matrix
	$(C, c)$. If $Cx\geq c$ has an integer solution, then it has
	a solution $x\in \setZ^d$ such that $|x(i)|\leq (d+1)\cdot m$ for
	all $i\in [1,d]$.
\end{lemC}

\begin{lem}[\textbf{Systems without divisibility constraints}]
\label{lem:nodivsmallsolutions}
\hfill\\
Every satisfiable $d$-dim linear system $S$ without divisibility constraints has a solution $x \in \mathbb{Z}^d$ such that 
\[
    \nnorm{x} \leq (d+1)\,!\cdot(\nnorm{S}+1)^d.
\]
\end{lem}
\begin{proof}
Let $S$ be a satisfiable $d$-dim linear system, expressed as a
	propositional formula whose atoms are equality and inequality
	constraints (with no divisibility constraints).
Each equality constraint
	$\scalar{\alpha}{\bvar{x}}=b$ in $S$ can be replaced by the
	conjunction
	$\scalar{\alpha}{\bvar{x}}\geq b\,\wedge\,
	\scalar{-\alpha}{\bvar{x}}\geq -b$ without affecting
	$\sem{S}$ or $\nnorm{S}$. 

	We may thus assume,
	without loss of generality, that the atoms in $S$ consist only
	of inequality constraints.
By
expressing $S$ in disjunctive normal form, we obtain
	linear systems $S_1,S_2,\ldots,S_k$, each given as a
	conjunction of inequalities and negations of
	inequalities, where
	\[\sem{S}=\bigcup_{j=1}^k\sem{S_j}\text{, and }
	\nnorm{S_j}\leq \nnorm{S}\text{ for all }j\in[1,k].\]
Since $\sem{S}\neq\emptyset$, it follows that $\sem{S_j}\neq\emptyset$ for some $j\in[1,k]$.
It therefore suffices to establish the claim for the case where  
$S$ is a~conjunction of inequality constraints and their negations; we
	will assume this henceforth.

Let $S'$ be the linear system obtained from $S$ by replacing
	each subformula $\neg(\scalar{\alpha}{\bvar{x}}\geq b)$ with
	the
	inequality constraint $\scalar{-\alpha}{\bvar{x}}\geq -b+1$.
Then
	\begin{equation}\label{eq:semSprime}\sem{S'}=\sem{S}=\{x\in\setZ^d\mid
	Cx\geq c\},\end{equation}
where $C$ and $c$ are some $n\times d$ and $n\times 1$ integer
	matrices, respectively, with
	\[
		\max\{\nnorm{C},\nnorm{c}\}=\nnorm{S'}\leq\nnorm{S}+1.\]
        Let $r$ be the rank of $C$, and let $m$ be an upper bound on
	the absolute values of the $(r-1)\times (r-1)$ and $r\times r$
	subdeterminants of the $n\times (d+1)$ augmented matrix $(C, c)$.
Since  $r\leq d$, it follows that
	\[m\leq d\,!\cdot (\nnorm{S}+1)^d.\]
	Applying Lemma~\ref{lem:Gathen} to~\eqref{eq:semSprime}, we
conclude that $S'$ admits a solution $x$ such that 
	\[\nnorm{x}\leq (d+1)\cdot m\leq (d+1)\,!\cdot
	(\nnorm{S}+1)^d.\qedhere\]
\end{proof}


\noindent
\textbf{Proof of Theorem~\ref{thm:smallsolutions}}. 
Let $S$ be a satisfiable $d$-dim linear system, with a~variable vector
$\bvar{x}$ ranging over $\setZ^d$. We set
$\ell=\mlcm{S}$ and fix a solution $\sol$ of $S$.
Let 
\begin{center}
$b\in \{0,1,\dots,\ell{-}1\}^d$ be the vector of remainders of
	$\sol$ modulo $\ell$;
\end{center}
specifically, we define $b(i)=(\sol(i)\bmod\ell)$ for all $i\in[1,d]$. 

We aim to construct a $d$-dim linear system $S_b$, with a variable
vector
$\bvar{y}$ and without divisibility constraints, such that 
\begin{equation}\label{eq:semSb}
\sem{S_b}=\{y\in \setZ^d\mid b+\ell y\in\sem{S}\}.
\end{equation}
We construct $S_b$ from $S$ by replacing each constraint as follows:
  \begin{itemize}
  \item $\scalar{\alpha}{\bvar{x}}=c$ is replaced by 
	  $\scalar{\alpha}{\bvar{y}}=\frac{c-\scalar{\alpha}{b}}{\ell}$
		  if $\ell$ divides $c-\scalar{\alpha}{b}$, and by
		  \emph{false} otherwise;
  \item $\scalar{\alpha}{\bvar{x}}\geq c$ is replaced by
	  $\scalar{\alpha}{\bvar{y}}\geq  \left
		  \lceil{\frac{c-\scalar{\alpha}{b}}{\ell}}\right
		  \rceil$;
  \item $\scalar{\alpha}{\bvar{x}}\equiv c\,(\bmod\, m)$ is replaced
	  by \emph{true} if
		  $\scalar{\alpha}{b}\equiv c\,(\bmod\, m)$ and by
		  \emph{false}
		  otherwise.
  \end{itemize}

% Fix Indentation
\noindent
It is straightforward to verify that $\sem{S_b}$ satisfies the 
condition~\eqref{eq:semSb}, by noting that $y\in\sem{\text{new
constraint}}$ if and only if
$x\in\sem{\text{old constraint}}$ for $x=b+\ell y$
(using the fact that $\ell$ is a multiple of each $m$ appearing in the
divisibility constraints).
Moreover, $\sem{S_b}\neq\emptyset$ because the vector $y$ defined by 
$y(i)=\frac{\sol(i)-b(i)}{\ell}$ for each $i\in[1,d]$ is a solution.


Note that if $\scalar{\alpha}{\bvar{x}}\sim c$ is a constraint of
$S$, where $\sim$ is the equality or the inequality, then
the values $|\alpha(i)|$ for $i\in[1,d]$ and $|c|$ are bounded by
$\nnorm{S}$. It follows that
$|\scalar{\alpha}{b}|\leq d\cdot\nnorm{S}\cdot(\ell-1)$, and 
 $|\frac{c-\scalar{\alpha}{b}}{\ell}|\leq
d\cdot\nnorm{S}$. We have thus shown that 
\[\nnorm{S_b}\leq d\cdot\nnorm{S}.\] 
Since $S_b$ is satisfiable and contains no divisibility constraints, 
Lemma~\ref{lem:nodivsmallsolutions} implies that $S_b$ has a solution
$y\in\setZ^d$ such that 
\[\nnorm{y}\leq (d+1)\,!\cdot(\nnorm{S_b}+1)^d.\]
Hence $x=b+\ell y$
is a solution of $S$, satisfying
\[\nnorm{x}\leq (\ell-1) +\ell\cdot (d+1)\,!\cdot(\nnorm{S_b}+1)^d\leq
	\ell\cdot\left(1+ (d+1)\,!\cdot(d\cdot \nnorm{S}+1)^d\right).\tag*{\qed}\]
